\chapter{Hodge理论与Gauss--Bonnet定理}
\label{ch:14-hodge-gauss-bonnet}
\ChapterPlan{周4.6至4.7（本地页码待核对）；梁上册第二版5.6及习题11--13、15；Donaldson椭圆讲义与Andrews第20讲补充}{Hodge及曲面Gauss--Bonnet证明初稿；8题详解与3幅原创图，尚未独立数学审阅；高维整体公式仅陈述}

微分形式的上同调不依赖度量，而度量提供了“同一个类中哪一个代表最合适”的问题。Hodge理论给出的答案是调和代表；存在性不能只靠形式上的正交分解，而须证明弱方程可解及其解光滑。另一方面，曲面的曲率积分由局部几何量拼成一个整数，这由Gauss--Bonnet解释。

本章的Hodge定理针对光滑、紧致、定向、无边界的有限维黎曼流形；星算子的逐点定义和紧支分部积分可在更一般的定向流形上使用。分析部分以Lebesgue积分及$L^2$的完备性为先修，补建所需的Hilbert变分引理、Sobolev紧性和弱解正则性，不调用尚未证明的谱分解。曲面部分以明确的有限光滑三角剖分为拓扑输入，证明边界/转角公式和求和；一般光滑三角剖分存在定理及其与奇异同调Euler数的比较不在本章证明。高维Chern--Gauss--Bonnet的整体等式另标为外引陈述，不参与前面证明。

\Needspace{17\baselineskip}
\section{Hodge星算子与余微分}
\subsection*{外幂内积与定向}
在$n$维定向内积空间中，协向量内积由逆度量给出。约定
\[
 \langle\alpha_1\wedge\cdots\wedge\alpha_k,
 \beta_1\wedge\cdots\wedge\beta_k\rangle
 =\det(\langle\alpha_i,\beta_j\rangle).
\]
取正向正交归一余标架$\theta^1,\ldots,\theta^n$，严格递增的$\theta^I$构成$\Lambda^kT^*M$的正交归一基。若用全指标分量写$\alpha=\alpha_{i_1\cdots i_k}\theta^{i_1}\wedge\cdots\wedge\theta^{i_k}/k!$，则$|\alpha|^2=\sum_{i_1,\ldots,i_k}\alpha_{i_1\cdots i_k}^2/k!$，不能与递增指标求和混用。

\begin{definition}[Hodge星]\label{def:hodge-star}
星算子$*:\Lambda^kT^*M\to\Lambda^{n-k}T^*M$由
\[
 \alpha\wedge *\beta=\langle\alpha,\beta\rangle\mathrm{vol}_g
 \quad\text{对所有$\alpha,\beta\in\Lambda^kT^*M$}
\]
唯一确定。具体地，$*\theta^I=\varepsilon(I,I^c)\theta^{I^c}$，符号取使$\theta^I\wedge *\theta^I=\mathrm{vol}_g$者。
\end{definition}
外积配对非退化给唯一性；补指标公式给存在性和光滑性，因此与正向正交标架的选择无关。移动补指标跨过$k(n-k)$次，得到
\[
 **|_{\Omega^k}=(-1)^{k(n-k)}I,\qquad
 *1=\mathrm{vol}_g,\qquad *\mathrm{vol}_g=1,
 \qquad |*\alpha|=|\alpha|.
\]
反转定向使$*$与体积形式同时变号。若$\widetilde g=e^{2u}g$，$k$形式的内积乘$e^{-2ku}$、体积乘$e^{nu}$，故
\[
 *_{\widetilde g}|_{\Omega^k}=e^{(n-2k)u}*_g.
\]
在中间次数$n=2k$星算子共形不变，但一般次数并不如此。

\begin{figure}[htbp]
\centering\DGDiagram{hodge-star-plane}

\caption{定向欧氏平面中$*\dd x=\dd y$、$*\dd y=-\dd x$，因此一次形式上$**=-I$。图经内积把协向量识别为箭头，表示星算子而非切向量的坐标换基。}
\label{fig:hodge-star-plane}
\end{figure}

\begin{proposition}[星算子与Levi--Civita导数]\label{prop:parallel-hodge-star}
Levi--Civita联络满足$\nabla\mathrm{vol}_g=0$和$\nabla_X(*\alpha)=*(\nabla_X\alpha)$。
\end{proposition}
\begin{proof}
局部正交标架的联络矩阵斜对称。顶外幂的诱导联络系数是其迹的负号，故为零，得到体积平行。外幂内积也因度量相容而平行。对定义中的$\alpha\wedge *\beta$作协变导数，扣除对$\alpha,\beta$求导的两项，得到$\alpha\wedge(\nabla_X(*\beta)-*\nabla_X\beta)=0$对所有$\alpha$成立；非退化性给结论。这里不是说坐标中$*$的系数为常数，而是联络项恰补偿其变化。
\end{proof}

\subsection*{形式伴随及两个Laplace符号}
对至少一个紧支的同次数实形式定义$(\alpha,\beta)=\int_M\alpha\wedge *\beta$。在无边界情形，这等于正定的$L^2$内积。定义
\[
 \delta|_{\Omega^k}=(-1)^{n(k+1)+1}*\dd*,\qquad \delta|_{\Omega^0}=0.
\]
\begin{proposition}[余微分与能量恒等式]\label{prop:codifferential-energy}
在无边界流形上，至少一个紧支时有$(\dd\alpha,\beta)=(\alpha,\delta\beta)$。并且$\delta^2=0$，Hodge Laplace算子
\[
 \Delta_H=\dd\delta+\delta\dd
\]
满足$\Delta_H\dd=\dd\Delta_H$、$\Delta_H\delta=\delta\Delta_H$及
\[
 (\Delta_H\alpha,\beta)=(\dd\alpha,\dd\beta)+(\delta\alpha,\delta\beta).
\]
在函数上$\Delta_Hf=-\Delta_gf$，其中第11章$\Delta_g=\operatorname{div}\operatorname{grad}$。
\end{proposition}
\begin{proof}
令$\deg\alpha=k-1$、$\deg\beta=k$。Stokes给
\[
 0=\int\dd(\alpha\wedge *\beta)
  =\int\dd\alpha\wedge *\beta+(-1)^{k-1}\int\alpha\wedge\dd*\beta.
\]
用$**$公式化简可得$*\delta\beta=(-1)^k\dd*\beta$，所以正好得到伴随式。将两次$\delta$中的中间$**$合并，$\dd^2=0$给$\delta^2=0$；两条交换关系分别展开即得。伴随式应用两次得到能量式。

在一点取联络系数为零的正交标架，外微分与余微分分别为
\[
 \dd\alpha=\sum_i\theta^i\wedge\nabla_{e_i}\alpha,
 \qquad\delta\alpha=-\sum_i\iota_{e_i}\nabla_{e_i}\alpha.
\]
第一式由无挠及外微分坐标式得到；第二式也可在此标架逐分量分部积分求伴随。于是$\delta\dd f=-\operatorname{tr}(\nabla\dd f)=-\Delta_gf$。若有边界，Stokes会留下边界项，本命题的伴随结论不能在没有边界条件时直接使用。
\end{proof}

在紧致无边界情形，$(\Delta_H\alpha,\alpha)=\|\dd\alpha\|^2+\|\delta\alpha\|^2\geq0$。称$\Delta_H\alpha=0$的光滑形式为调和形式，记空间为$\mathcal H^k$；此时调和等价于同时闭且余闭。这个积分论证还没有证明每个上同调类有调和代表。

\subsection*{主符号为什么正定}
记$\varepsilon_\xi(\alpha)=\xi\wedge\alpha$，$\iota_\xi=\iota_{\xi^\sharp}$。插入的分次乘法式给
\[
 \iota_\xi\varepsilon_\xi+\varepsilon_\xi\iota_\xi=|\xi|^2I.
\]
在最高阶部分用$\partial_j\mapsto i\xi_j$定义主符号，则
\[
 \sigma_{\dd}(\xi)=i\varepsilon_\xi,\qquad
 \sigma_\delta(\xi)=-i\iota_\xi,\qquad
 \sigma_{\Delta_H}(\xi)=|\xi|^2I.
\]
所以$\Delta_H$椭圆，局部最高阶项为$-g^{ij}\partial_i\partial_j$乘纤维恒等矩阵。这里的正号对应非负$\Delta_H$；若漏掉Fourier符号中的$i$，便会与第11章的$\Delta_g$混淆。


\Needspace{17\baselineskip}
\section{Sobolev工具与椭圆方程}
\subsection*{分析先修与Hilbert空间工具}
以下只用实数或有限个复数分量；Fourier计算暂用复系数，最后取实部。本节接受Lebesgue积分、$L^2$完备性及有限维光滑换元作为分析先修。光滑函数的密度可由截断、简单函数的矩形逼近和卷积得到；平移在$L^2$中连续，先对矩形特征函数验证，再用稠密性推广。后面的紧性和弱方程可解性均在这里补证。

\begin{lemma}[正交投影、Riesz与对称变分]\label{lem:hilbert-variational}
设$H$为Hilbert空间。闭子空间有唯一正交投影，连续线性泛函$\ell$可唯一写为$\ell(v)=(v,w)$。若对称双线性型$a$满足
\[
 |a(u,v)|\leq C\|u\|_H\|v\|_H,\qquad
 a(u,u)\geq c\|u\|_H^2\quad(c>0),
\]
则对每个连续$\ell$，存在唯一$u\in H$满足$a(u,v)=\ell(v)$对所有$v\in H$成立，且$\|u\|_H\leq\|\ell\|/c$。
\end{lemma}
\begin{proof}
对给定$x$取闭子空间中的距离极小化序列$y_j$。平行四边形恒等式与中点仍在子空间，给
\[
 \|y_j-y_l\|^2\leq
 2\|x-y_j\|^2+2\|x-y_l\|^2-4\inf_y\|x-y\|^2\longrightarrow0.
\]
完备性给极小元；沿子空间任意方向微扰并令距离平方的一次项为零，知余量正交，亦给唯一性。

若$\ell\ne0$，在闭仿射超平面$\ell(v)=1$上以同一中点论证取得最小范数元$z$，并有$z\perp\ker\ell$。分解$v=\ell(v)z+(v-\ell(v)z)$得$\ell(v)=(v,z)/\|z\|^2$；零泛函另计。最后$a$本身定义与原范数等价的Hilbert内积，故在该内积下对$\ell$应用刚证明的表示，得到唯一$u$。取$v=u$给所列估计。这是本章需要的对称Lax--Milgram情形，无须先调用一般Fredholm理论。
\end{proof}

可分Hilbert空间中的有界序列有弱收敛子列：在一组可数正交基的各系数上逐次取子列，再对角化；系数平方和由有界性控制，Riesz表示给极限向量，先对有限线性组合再用一致有界推广。后面出现的Sobolev空间均可分，因为局部Fourier多项式构成可数稠密集合。

\subsection*{在平坦环面上建立Sobolev工具}
取$\mathbb T^n=(\R/2\pi\mathbb Z)^n$，将测度归一为$(2\pi)^{-n}\dd x$。指数函数$e^{i\nu\cdot x}$正交归一。它们在$L^2$中完备，可用Fejer核说明：一维
\[
 K_N(t)=\frac1N\left|\sum_{j=0}^{N-1}e^{ijt}\right|^2\geq0,
 \quad\frac1{2\pi}\int K_N=1,
 \quad\int_{|t|\geq\epsilon}K_N\longrightarrow0.
\]
最后一式由几何级数及$|\sin(t/2)|\geq\sin(\epsilon/2)$得到。在多维取乘积核。利用$L^2$平移连续性和三角不等式，$K_N*u\to u$于$L^2$；卷积是有限Fourier多项式，故多项式稠密，正交投影的极限给Parseval恒等式。

对非负整数$s$定义
\[
 \|u\|_{H^s(\mathbb T^n)}^2
 =\sum_{\nu\in\mathbb Z^n}(1+|\nu|^2)^s|\widehat u(\nu)|^2.
\]
逐项求导与Parseval说明它等价于所有$|\gamma|\leq s$阶弱偏导的$L^2$范数平方和；弱偏导按测试函数分部积分定义。Fourier截断同时逼近各阶导数，故光滑函数在$H^s$中稠密，且这些加权序列空间完备。

\begin{lemma}[Sobolev连续性与Rellich紧性]\label{lem:sobolev-rellich}
若整数$s>q+n/2$，则$H^s(\mathbb T^n)$连续嵌入$C^q$。$H^1$中的有界集合在$L^2$中相对紧。
\end{lemma}
\begin{proof}
对$|\gamma|\leq q$，Cauchy--Schwarz给
\[
 \sum_\nu |\nu^\gamma\widehat u(\nu)|
 \leq\|u\|_{H^s}
 \left(\sum_\nu(1+|\nu|^2)^{q-s}\right)^{1/2}<\infty.
\]
级数收敛由按整数壳计数，壳中格点数为$O(R^{n-1})$。故各阶Fourier导数一致收敛，给连续性及范数控制。

若$\|u_j\|_{H^1}\leq C$，截去$|\nu|>R$的尾部，其$L^2$范数平方至多$C^2/(1+R^2)$。低频只在有限维有界集合中，存在收敛子列。对$R=1,2,\ldots$对角化，尾部一致小使最终子列在$L^2$中Cauchy，完备性给强收敛。这一尾部估计是紧性的来源，不能只由有界性推出强收敛。
\end{proof}

\subsection*{搬到紧流形的向量丛}
在有限坐标覆盖上取紧支单位分解$\chi_a$及局部丛标架，把每个$\chi_au$的系数支持缩放到$(-\pi,\pi)^n$内部，再延零为环面函数。定义$\|u\|_{H^s}^2$为这些系数$H^s$范数的总和；对光滑截面取完备化。坐标变换的链式法则、紧支集上Jacobian上下界及标架变换的有界导数，说明不同有限选择给等价范数。

弱导数的局部描述与上述完备化一致：一个方向由收敛后对测试函数分部积分得到；另一方向在坐标中先卷积平滑，再用单位分解拼接，卷积对每阶弱导数收敛，得到光滑逼近。因而$H^{s+1}\to H^s$连续，一阶光滑系数算子从$H^{s+1}$连续到$H^s$。由有限覆盖逐片应用上引理，有
\[
 H^s(M;E)\hookrightarrow C^q(M;E)\quad(s>q+n/2),
 \qquad H^1(M;E)\hookrightarrow L^2(M;E)\quad\text{紧}.
\]
对后一个结论，在有限片上依次取强收敛子列即可。此处流形无边界；没有在坐标边界上延零任意$H^1$函数，延零的是支集严格位于卡内的$\chi_au$。

\subsection*{一阶能量估计}
以下流形紧致且无边界，次数$k$固定。写
\[
 Q(u,v)=(\dd u,\dd v)+(\delta u,\delta v),\qquad
 Q(u,u)=\|\dd u\|^2+\|\delta u\|^2.
\]
一阶算子的连续性把$Q$延拓到$H^1$形式。

\begin{lemma}[Garding估计]\label{lem:hodge-garding}
存在仅依赖流形、度量和次数的$C$，对所有$u\in H^1(\Lambda^kT^*M)$有
\[
 \|u\|_{H^1}^2\leq C\bigl(Q(u,u)+\|u\|_{L^2}^2\bigr).
\]
\end{lemma}
\begin{proof}
先对支集在小坐标球中的光滑形式证明。在球中心将度量及纤维内积化为标准内积，令$D=\dd+\delta$，冻结其一阶系数为常系数算子$D_0$。主符号恒等式给
\[
 |\varepsilon_\xi a|^2+|\iota_\xi a|^2=|\xi|^2|a|^2,
 \qquad \|\partial u\|_{L^2}=\|D_0u\|_{L^2}.
\]
第二式可在包含支集的环面上逐Fourier系数验证。由于一阶系数连续，缩小球后
\[
 \|(D-D_0)u\|_{L^2}
 \leq\epsilon\|\partial u\|_{L^2}+C_\epsilon\|u\|_{L^2}.
\]
取$\epsilon<1/2$吸收到左边，便有局部估计。这里把$Du$视为$k-1$与$k+1$次两部分的直和，因此$\|Du\|^2=Q(u,u)$，没有不同次数相消的问题。

在有限覆盖上对$\chi_au$应用局部估计。交换子
\[
 [\dd,\chi_a]u=\dd\chi_a\wedge u,\qquad
 [\delta,\chi_a]u=-\iota_{\operatorname{grad}\chi_a}u
\]
均为零阶，平方和由$C\|u\|^2$控制。有限个局部估计相加，得到全局式。最后用光滑稠密性推广到$H^1$。低阶项解释了为何估计中必须保留$L^2$范数。
\end{proof}

\subsection*{弱解的二阶正则性}
称$u\in H^1$为$\Delta_Hu=f\in L^2$的弱解，是指$Q(u,v)=(f,v)$对所有光滑测试形式$v$成立；由稠密性也对所有$H^1$测试成立。

\Needspace{18\baselineskip}
\begin{lemma}[从弱解到光滑解]\label{lem:hodge-weak-regularity}
局部$H^1$弱解满足：若$f\in H^s_{\mathrm{loc}}$、$s\geq0$为整数，则$u\in H^{s+2}_{\mathrm{loc}}$。在紧致无边界流形上，存在$C_s$使
\[
 \|u\|_{H^{s+2}}\leq C_s\bigl(\|f\|_{H^s}+\|u\|_{L^2}\bigr).
\]
特别$f$光滑时$u$光滑。该结论针对光滑系数的Hodge算子，不声称任意粗糙系数系统具有同样正则性。
\end{lemma}
\begin{proof}
在局部光滑正交丛标架中，主部为标量矩阵；令$\rho=\sqrt{\det g}$，将方程乘$\rho$，可写为
\[
 -\partial_i(a^{ij}\partial_j u)+C^j\partial_ju+C_0u=\rho f,
 \qquad a^{ij}=\rho g^{ij}I.
\]
所有矩阵系数光滑，紧子球上$a$一致正定。此形式来自前述主符号，其余二阶项既已扣除，差只能为一阶与零阶算子；局部弱形式可由测试函数分部积分得到。

下面给出从$H^1$到$H^2$的差商论证。固定一个坐标方向，记$D_hw(x)=(w(x+he_l)-w(x))/h$。取内部截断函数$\eta$，$|h|$小到所有平移支集仍在坐标球内，以
\[
 v=-D_{-h}(\eta^2D_hu)
\]
测试方程。对固定$h$，此$v$为紧支$H^1$函数，可用稠密性作测试。平移换元给$\int pD_{-h}q=-\int(D_hp)q$。主项因此含
\[
 \int\eta^2 a^{ij}(x+he_l)
       \partial_jD_hu\cdot\partial_iD_hu
 \geq c\|\eta\,\partial D_hu\|_{L^2}^2.
\]
其余项来自$\partial\eta$、$D_ha^{ij}$及低阶系数。利用
\[
 D_h(a\partial u)=a(x+he_l)\partial D_hu+(D_ha)\partial u,
 \qquad \|D_hw\|_{L^2}\leq\|\partial_lw\|_{L^2},
\]
它们的绝对值至多
$C\|u\|_{H^1}(\|\eta\partial D_hu\|_{L^2}+\|u\|_{H^1})$。
右端无需对$f$求导，因为
\[
 \left|\int\rho fD_{-h}(\eta^2D_hu)\right|
 \leq C\|f\|_{L^2}
       \bigl(\|\eta\partial D_hu\|_{L^2}+\|u\|_{H^1}\bigr).
\]
这里再次用差商的导数界作用于$\eta^2D_hu$。用$ab\leq\epsilon a^2+b^2/(4\epsilon)$吸收，得到与$h$无关的内球估计
\[
 \|\eta\partial D_hu\|_{L^2}
 \leq C(\|f\|_{L^2}+\|u\|_{H^1}).
\]
当$h\to0$时，二阶差商有弱收敛子列；对测试函数分部积分表明弱极限正是二阶分布导数。逐方向得到$u\in H^2$及局部估计。上面的差商导数界由$w(x+h)-w(x)=\int_0^h\partial_lw(x+t)\dd t$、Cauchy--Schwarz和密度得到，故也适用于$H^1$函数。

归纳高阶：若已知$u\in H^{s+1}_{\mathrm{loc}}$且$f\in H^s_{\mathrm{loc}}$，对$|\gamma|=s$在分布意义求导，
\[
 L(\partial^\gamma u)=\partial^\gamma(\rho f)+[L,\partial^\gamma]u.
\]
交换子至多含$s+1$阶的$u$导数：每个未抵消的最高阶项至少有一次导数落在光滑系数上，故右端在$L^2$。对$\partial^\gamma u\in H^1$再用刚证的二阶结论，得到$u\in H^{s+2}$。取嵌套内球并有限次应用，给相应局部范数估计。

紧流形用有限覆盖相加，先得到右端含$\|u\|_{H^1}$的全局估计。对弱方程取$v=u$并用Garding式，
\[
 \|u\|_{H^1}^2\leq C((f,u)+\|u\|^2)
 \leq C'(\|f\|^2+\|u\|^2),
\]
于是得到所列估计。若$f$光滑，归纳得$u$属于全部整数阶Sobolev空间，再由嵌入定理得全部连续阶导数，即为光滑解。
\end{proof}


\Needspace{17\baselineskip}
\section{Poincare对偶与调和形式}
\subsection*{先排除零能量方向，再解弱方程}
考虑$\mathcal N=\{u\in H^1:\dd u=0,\delta u=0\}$。这些等式在$L^2$或分布意义下相同。对$u\in\mathcal N$有$Q(u,v)=0$，正则性引理使$u$光滑，故$\mathcal N=\mathcal H^k$。

\begin{lemma}[有限维核与核外强制性]\label{lem:hodge-coercivity}
$\mathcal H^k$有限维。存在$c>0$，使所有$u\in H^1$且$u\perp\mathcal H^k$满足
\[
 Q(u,u)\geq c\|u\|_{H^1}^2.
\]
\end{lemma}
\begin{proof}
若核无限维，取其中$L^2$正交归一序列。Garding估计使其$H^1$有界，Rellich给$L^2$强收敛子列；但任意不同两项距离为$\sqrt2$，矛盾。

先证明核外$\|u\|_{L^2}^2\leq C Q(u,u)$。若否，存在$\|u_j\|_{L^2}=1$、$u_j\perp\mathcal H^k$且$Q(u_j,u_j)\to0$。Garding使$H^1$有界，取子列在$H^1$弱收敛、$L^2$强收敛到$u$。一阶算子连续给$\dd u=\delta u=0$，所以$u\in\mathcal H^k$；强收敛保留范数1及与核正交，矛盾。将此界代入Garding，便得$H^1$强制性。
\end{proof}

\Needspace{19\baselineskip}
\begin{theorem}[Hodge分解与调和代表]\label{thm:hodge-decomposition}
设$(M,g)$光滑、紧致、定向、无边界，则每个次数有$L^2$正交直和
\[
 \Omega^k(M)=\mathcal H^k\oplus\dd\Omega^{k-1}(M)
                        \oplus\delta\Omega^{k+1}(M).
\]
约定越界次数的空间为零。每个de Rham上同调类含唯一调和代表，因此$\mathcal H^k\simeq H^k_{\mathrm{dR}}(M;\R)$，特别上同调有限维。
\end{theorem}
\begin{proof}
令$P$为有限维光滑核的$L^2$正交投影，取光滑$\alpha$并置$f=\alpha-P\alpha$。空间
\[
 V=H^1(\Lambda^kT^*M)\cap(\mathcal H^k)^\perp
\]
在$H^1$中闭，因此完备。$Q$在$V$连续、对称且强制，泛函$v\mapsto(f,v)$连续；变分引理给唯一$u\in V$使$Q(u,v)=(f,v)$对$v\in V$成立。任意$H^1$测试形式可减去其核投影，而$Q(u,h)=0=(f,h)$对核成立，故此弱方程对全部测试成立。正则性使$u$光滑，且
\[
 \alpha=P\alpha+\Delta_Hu=P\alpha+\dd\delta u+\delta\dd u.
\]
调和形式因闭且余闭而与两个像正交；两像也正交，因为
\[
 (\dd\beta,\delta\gamma)=(\dd\dd\beta,\gamma)=0.
\]
所以这是正交直和，各个分量唯一；势$\beta,\gamma$本身一般不唯一。

若$\dd\alpha=0$，对分解中的余恰当项配对，
\[
 (\alpha,\delta\dd u)=(\dd\alpha,\dd u)=0.
\]
由正交性得$\|\delta\dd u\|^2=0$，因此$\alpha=P\alpha+\dd\delta u$，与调和代表同类。若调和$h$又恰当，$h=\dd\beta$，则$\|h\|^2=(h,\dd\beta)=(\delta h,\beta)=0$，故$h=0$。这同时给唯一性和同构，而不是仅仅比较两个空间的维数。
\end{proof}

记上述核外解算子为Green算子$G$，规定$G\alpha\perp\mathcal H^k$且$GP=0$。则在光滑形式上
\[
 \Delta_HG=G\Delta_H=I-P,\qquad
 \alpha=P\alpha+\dd\delta G\alpha+\delta\dd G\alpha.
\]
第一式左半已证明。右半用自伴性得$\Delta_H\alpha\perp\mathcal H^k$，而$\alpha-P\alpha$是相应核外解，再用唯一性。变分估计和正则性还给$\|Gf\|_{H^{s+2}}\leq C_s\|f\|_{H^s}$；这并不预设有一个逐点Green核。

\begin{figure}[htbp]
\centering\DGDiagram{hodge-decomposition}

\caption{Hodge分解的有限维投影示意：三个分量位于调和、恰当、余恰当子空间，正交性是$L^2$积分意义的，不是逐点正交；图中三轴是空间投影。调和项由有限维投影给出，另外两项需先解椭圆方程。}
\label{fig:hodge-decomposition}
\end{figure}

\subsection*{对偶配对及其度量依赖}
\begin{theorem}[紧定向流形的de Rham对偶]\label{thm:hodge-poincare-duality}
同样假设下，积分配对
\[
 H^k_{\mathrm{dR}}(M)\times H^{n-k}_{\mathrm{dR}}(M)\longrightarrow\R,
 \qquad([\alpha],[\beta])\longmapsto\int_M\alpha\wedge\beta
\]
非退化，因而$H^{n-k}_{\mathrm{dR}}\simeq(H^k_{\mathrm{dR}})^*$。这是普通de Rham上同调的对偶，不是尚未建立的非紧/紧支版本。
\end{theorem}
\begin{proof}
闭形式与恰当形式的外积积分由Stokes为零，故配对与代表选择无关。若$h\in\mathcal H^k$，由$\delta$的星公式可知$*h$同时闭且余闭，因此$*h\in\mathcal H^{n-k}$。若$[\alpha]\ne0$，其唯一调和代表$h\ne0$，取$\beta=*h$，则
\[
 \int_Mh\wedge *h=\|h\|_{L^2}^2>0.
\]
交换两个因子同理，得到双侧非退化。两空间有限维，所以到对方对偶空间的单射在两边维数不等式下均为同构。
\end{proof}

积分配对本身只用定向及上同调；通过调和代表构造$[h]\mapsto[*h]$时使用了度量，不能说星算子在任意闭形式上直接给出同一个与度量无关的上同调映射。零次调和形式是逐连通分支常数，顶次由星对应各分支体积形式。若$M$连通，顶次积分给$H^n\simeq\R$。

对平坦$\mathbb T^n$，$\Delta_H$逐系数为$-\sum_i\partial_i^2$，Fourier零频以外均无核，所以调和形式恰为常系数形式，$\dim H^k=\binom nk$。球面的上同调已在第9章计算，故其调和形式只在0与$n$次非零；结论不依赖使用圆度量。

\subsection*{假设不能删去}
在非紧$\R$上，函数$x$满足$\Delta_Hx=0$但$\dd x\ne0$；没有有限能量积分恒等式可用。一次形式$\dd x$既调和又恰当，故普通上同调中的调和代表不唯一。在有边界的区间上，同样的形式也给反例，除非另配绝对或相对边界条件。非定向流形上的普通全局星及积分配对则须改用定向线丛或密度版本。本章不把这些修订包含在刚证明的定理中。


\Needspace{17\baselineskip}
\section{曲面Gauss--Bonnet及高维推广}
\subsection*{联络一形式、测地曲率与转角}
设$S$为定向黎曼曲面，$J$为切平面中的正向直角旋转。局部取正向正交标架$(e_1,e_2)$，令
\[
 a(X)=g(\nabla_Xe_1,e_2)=\omega^2{}_1(X).
\]
第13章的曲率矩阵约定给$\Omega^1{}_2=K\,\mathrm{vol}_g$，而$\omega^1{}_2=-a$，所以
\[
 \dd a=-K\,\mathrm{vol}_g.
\]
沿单位速正则曲线，写$T=\cos\vartheta\,e_1+\sin\vartheta\,e_2$，则
\[
 k_g:=g(\nabla_TT,JT)=\frac{\dd\vartheta}{\dd s}+a(T).
\]
直接对$T$求导、用$\nabla e_1=ae_2$和$\nabla e_2=-ae_1$即可验证。正向边界以区域在左侧为准，故$JT$指向区域内部。若标架旋转$\phi$，有$a'=a+\dd\phi$、$\vartheta'=\vartheta-\phi$，所以$k_g$不变。

在边界角点，设区域内角为$\beta\in(0,2\pi)$，定义带号外转角$\epsilon=\pi-\beta$。凸角外转角为正，凹角为负；光滑边界上人为插入的顶点有$\beta=\pi$、$\epsilon=0$。

\Needspace{18\baselineskip}
\begin{lemma}[一个坐标三角片上的Gauss--Bonnet]\label{lem:local-gauss-bonnet}
设$T$为正向曲面卡内的闭三角片，它是平面闭三角形在某个邻域微分同胚下的像。沿三条光滑边取正向边界参数，则
\[
 \int_T K\,\mathrm{vol}_g+\int_{\partial T}k_g\,\dd s
        +\sum_{p\text{为角点}}\epsilon_p=2\pi.
\]
对同样由平面闭圆盘微分同胚参数化的光滑圆盘，公式没有角点项。
\end{lemma}
\begin{proof}
卡内可用Gram--Schmidt构造正向正交标架。在每条边上提升切向角$\vartheta$，到角点补上有号的外转角，使切方向成为一个连续绕行。此绕行的总角增量为$2\pi$，理由如下。

把切向量拉回参考三角形，并把选定标架拉回。切方向在标架中的分量由一个光滑$GL^+(2,\R)$矩阵$A(x)$作用于参考边界切方向得到。取三角形内点$p$，用$A_t(x)=A(p+t(x-p))$从$A$连续变到常矩阵$A(p)$。这些矩阵均可逆，边界切向量始终非零，角点处的正向插值也连续；因此角增量除以$2\pi$这个整数不变。常正行列式矩阵可在$GL^+(2,\R)$内连到恒等矩阵（Gram--Schmidt分解成旋转与正对角上三角矩阵，再分别变形），故总绕行等于参考三角形的总外转角$2\pi$。圆盘用相同的径向收缩和标准圆的切方向验证。

把$k_g=\vartheta'+a(T)$在各边积分并加角点项，得到
\[
 \int_{\partial T}k_g\,\dd s+\sum\epsilon_p
 =2\pi+\int_{\partial T}a
 =2\pi-\int_TK\,\mathrm{vol}_g,
\]
最后一步为Stokes。对于三角片的有限角点，可先在卡内截去半径$\epsilon$的小角并光滑化，应用第8章光滑边界Stokes；新增弧长为$O(\epsilon)$、删去面积为$O(\epsilon^2)$，而$a,\dd a$在闭片上有界，令$\epsilon\to0$即得所用的分段光滑公式。这只补证本节二维有限角点的情形，不预设一般角点流形理论。三条边无需是测地线，也没有默认任意闭曲线的绕数都是1。
\end{proof}

若三条边均为测地线，则角为$A,B,C$时
\[
 \int_T K\,\mathrm{vol}_g=A+B+C-\pi.
\]
正曲率球面出现角超，平坦时回到平面三角形内角和。
\begin{figure}[htbp]
\centering\DGDiagram{gauss-bonnet-octant}

\caption{半径$a$球面的第一卦限三角片，由三条大圆弧围成，每个内角为$\pi/2$。图为球面片的空间投影，三个直角由环境中相交大圆的切向量验证；面积为球面的八分之一$\pi a^2/2$，曲率积分恰为角超$\pi/2$。}
\label{fig:gauss-bonnet-octant}
\end{figure}

\subsection*{有限剖分上的全局求和}
本节明确拓扑输入：设紧曲面给定一个有限三角剖分，每个闭面具有上引理的坐标三角片参数，边界为边的并，角点列在顶点中，顶点处相邻片的切扇形恰覆盖整圆或边界内角。定义$\chi_{\mathcal T}=V-E+F$。曲面光滑三角剖分的一般存在性不在这里补证；下面的几何与组合论证对给定的这种剖分是完整的。

\Needspace{18\baselineskip}
\begin{theorem}[曲面Gauss--Bonnet，含边界与角点]\label{thm:surface-gauss-bonnet}
设$S$为具有上述有限兼容剖分的紧定向黎曼曲面，边界分段光滑，角点外转角为$\epsilon_p$。则
\[
 \int_S K\,\mathrm{vol}_g+\int_{\partial S}k_g\,\dd s
       +\sum_{p\in\partial S}\epsilon_p=2\pi\chi_{\mathcal T}.
\]
无边界时为$\int_SK\,\mathrm{vol}_g=2\pi(V-E+F)$。同一光滑曲面不同兼容剖分的$V-E+F$相同，记为$\chi(S)$；闭曲面的总曲率亦与度量无关。
\end{theorem}
\begin{proof}
先对每个面片应用局部公式并相加。内部边被相反方向经过两次。反向时$T\mapsto-T$、$JT\mapsto-JT$，$\nabla_{-T}(-T)=\nabla_TT$，故$k_g$变号，内部边积分相消。仅剩真实边界积分。

记内部/边界顶点数为$V_i,V_b$，边数为$E_i,E_b$。所有面片角数为$3F$，内部顶点周围内角和为$2\pi$，边界顶点$p$周围内角和为真实边界内角$\beta_p$。所以面片外转角的总和为
\[
 3\pi F-2\pi V_i-\sum_{p\in\partial S}\beta_p.
\]
面公式右端相加为$2\pi F$，移项并加上真实边界外转角$\sum_p(\pi-\beta_p)$，得到
\[
 \int_SK+\int_{\partial S}k_g+\sum_p\epsilon_p
 =2\pi V_i+\pi V_b-\pi F.
\]
此行积分中的测度沿用前式。剖分的计数关系为$3F=2E_i+E_b$，而每个边界分支闭合，故$E_b=V_b$。因此
\[
 2\pi(V-E+F)=2\pi(V_i+V_b-E_i-E_b+F)
             =2\pi V_i+\pi V_b-\pi F,
\]
即得结论；闭曲面取所有边界计数为零。

固定度量后，等式左端不依赖剖分，所以不同兼容剖分给相同整数。固定剖分而改变度量，闭曲面的右端不变，所以总曲率与度量无关。对任意紧光滑曲面无条件应用时须另用光滑三角剖分存在定理；把该整数与奇异同调或de Rham Betti数的交错和作一般比较，还需相应拓扑比较定理，不能仅由前面的Hodge对偶自动得到。
\end{proof}

\subsection*{两个符号检查与曲率障碍}
球面半径$a$有$K=a^{-2}$、面积$4\pi a^2$，故总曲率$4\pi$对应$\chi=2$。平坦环面有$K=0$，故$\chi=0$；任意兼容剖分可用于核对此整数。若给定曲面$\chi<0$，它不可能有处处非负曲率的度量；如果$\chi=0$且$K\geq0$，连续性及积分为零迫使$K\equiv0$。这些结论不包含曲面度量的存在性或整体分类。

在半径$R$的球面非极点标架卡中，$a=\cos\theta\,\dd\phi$（此式不延伸到极点），$\dd a=-\sin\theta\,\dd\theta\wedge\dd\phi=-K\mathrm{vol}$。北极帽边界$\theta=\theta_0$取$\phi$增加方向，则$k_g=\cot\theta_0/R$、$\dd s=R\sin\theta_0\dd\phi$，所以
\[
 \int_{\text{帽}}K\,\mathrm{vol}=2\pi(1-\cos\theta_0),\qquad
 \int_{\partial\text{帽}}k_g\,\dd s=2\pi\cos\theta_0.
\]
两者相加为$2\pi$。若直接在含极点的整个帽上把这个奇异标架的$a$用于Stokes，会漏掉$2\pi$；局部圆盘公式使用的是在整个圆盘光滑的标架。

\subsection*{高维Euler形式：可计算的部分与尚未证明的整体等式}
令$n=2m$，正向正交标架的曲率二形式矩阵为$\Omega^i{}_j$，简写$\Omega_{ij}$，其矩阵斜对称。定义
\[
 e_g=\frac1{(2\pi)^m2^m m!}
 \sum_{\sigma\in S_{2m}}\operatorname{sgn}(\sigma)
 \Omega_{\sigma(1)\sigma(2)}\wedge\cdots\wedge
 \Omega_{\sigma(2m-1)\sigma(2m)}.
\]
这就是$\operatorname{Pf}(\Omega)/(2\pi)^m$。为看清不变性，在辅助外代数中引入正向基$z_1,\ldots,z_{2m}$，令
\[
 \Xi=\frac12\sum_{i,j}\Omega_{ij}z_i\wedge z_j.
\]
$\Xi^m/m!$的顶次系数恰为上述Pfaffian。正交保向换基只使顶外幂乘行列式1，所以局部形式拼成全局$2m$形式。这里曲率二形式彼此交换，不应套用一次形式换序的负号。

二维时$e_g=\Omega_{12}/(2\pi)=K\mathrm{vol}/(2\pi)$。四维时直接展开得到
\[
 e_g=\frac1{4\pi^2}\bigl(\Omega_{12}\wedge\Omega_{34}
 -\Omega_{13}\wedge\Omega_{24}+\Omega_{14}\wedge\Omega_{23}\bigr).
\]
它是顶次形式，故自动闭；若将同一不变多项式用于更高维底空间上的定向偶秩丛，则闭性来自Bianchi，可留到第18章建立。对本章的球面曲率约定，二维归一化为正，不另插入未说明的整体负号。

\begin{quote}\small
\textbf{外引陈述，未在本章证明：}Chern--Gauss--Bonnet断言，对紧致、定向、无边界的$2m$维黎曼流形，$\int_Me_g=\chi(M)$。一般证明还需Euler类与零截面/指标的比较或完整的传递形式构造；本章没有建立这些工具。四维表达式及下列乘积例子只是局部代数和归一化检查，不作为一般定理的证明。此整体定理的证明留至示性类与指标专题。
\end{quote}
例如$S^2(a)\times S^2(b)$的混合曲率为零，只有两个因子曲率块，故$e_g=(4\pi^2a^2b^2)^{-1}\mathrm{vol}$，积分为4；具体推导见习题。奇数维时也不能把同样的Pfaffian表达式原样使用。


\section{Hodge与Gauss--Bonnet习题及详细解答}
14.1改述梁书第5章原题11（定理5-6-1及5-6-2），只取本章正定情形；14.2合并改述原题12、13；14.3改述原题15并增加删点反例。原题均在印刷131/PDF147，正文5.6在印刷121--124/PDF137--140。14.4--14.8为原创补充，不另造原书题号。

\Needspace{12\baselineskip}
\subsection*{习题14.1：两次星与三维向量运算}
改述梁书原题11的正定部分。证明$**\alpha=(-1)^{k(n-k)}\alpha$，并在定向欧氏$\R^3$把梯度、旋度、散度写成$\dd,*$的组合。解释四维二形式上的$**$与Lorentz版本为何不能直接混用。
\begin{proof}[解答]
取正向正交余标架，$*\theta^I=\varepsilon(I,I^c)\theta^{I^c}$。再取星，补指标移回原次序需$k(n-k)$次交换，故两次星的系数为$(-1)^{k(n-k)}$；对基线性推广即得。正定四维二形式满足$**=I$，有自对偶及反自对偶分解$\alpha_\pm=(\alpha\pm*\alpha)/2$。Lorentz内积的符号会额外改变两次星的结果，不能把正定的实特征值分解直接搬过去。

对向量场$A$令$\alpha=A^\flat$，逐分量有
\[
 (\operatorname{grad}f)^\flat=\dd f,\qquad
 (\operatorname{curl}A)^\flat=*\dd\alpha,\qquad
 \operatorname{div}A=*\dd*\alpha=-\delta\alpha.
\]
例如$\dd\alpha$的$\dd x\wedge\dd y$系数为$\partial_xA_y-\partial_yA_x$，取星后成为$\dd z$分量；其余循环同理。$*\alpha=A_x\dd y\wedge\dd z+A_y\dd z\wedge\dd x+A_z\dd x\wedge\dd y$，外微分后恰为散度乘体积，再取星得到散度。梯度是度量将$\dd f$升指标后的场，不把一形式本身叫切向量。
\end{proof}

\Needspace{12\baselineskip}
\subsection*{习题14.2：笛卡儿与球坐标中的星}
梁书原题12、13合并并扩展。验证$*\dd x=\dd y\wedge\dd z$及$*(\dd x\wedge\dd y\wedge\dd z)=1$，再对三维欧氏球坐标求$*\dd r$、$*\dd\theta$、$*\dd\phi$。说明坐标奇点不能用同一个公式强行代入。
\begin{proof}[解答]
笛卡儿正向正交余标架的补指标规则直接给前两式。球坐标中取$r>0$、$0<\theta<\pi$且避开经线接缝的卡，正向正交余标架为
\[
 e^1=\dd r,\qquad e^2=r\dd\theta,\qquad e^3=r\sin\theta\dd\phi.
\]
因此
\begin{align*}
 *\dd r&=r^2\sin\theta\,\dd\theta\wedge\dd\phi,\\
 *\dd\theta&=\sin\theta\,\dd\phi\wedge\dd r,\\
 *\dd\phi&=\frac1{\sin\theta}\,\dd r\wedge\dd\theta.
\end{align*}
例如$\dd\theta\wedge*\dd\theta=\sin\theta\,\dd r\wedge\dd\theta\wedge\dd\phi=|\dd\theta|^2\mathrm{vol}$，核对了次序及逆度量因子。极点处$\phi$不是光滑坐标，$\dd\phi$本身无此全局意义；含$1/\sin\theta$的表达式不表示星算子在流形上奇异。原点$r=0$亦须换卡。
\end{proof}

\Needspace{12\baselineskip}
\subsection*{习题14.3：向量势的拓扑限制}
改述梁书原题15并扩展。证明整个$\R^3$上光滑无旋场有标量势、无散场有向量势。给出删去原点后无散却没有全局向量势的例子。
\begin{proof}[解答]
若$\operatorname{curl}A=0$，上题说明$\dd(A^\flat)=0$；$\R^3$为星形域，第7章Poincare引理给$A^\flat=\dd f$，故$A=\operatorname{grad}f$。若$\operatorname{div}B=0$，则$\dd*(B^\flat)=0$，同理存在全局一形式$\beta$使$*B^\flat=\dd\beta$。三维二形式上$**=1$，所以$B^\flat=*\dd\beta$，即$B=\operatorname{curl}(\beta^\sharp)$。势并不唯一，添加梯度不改变旋度。

在$\R^3\setminus\{0\}$取$B=x/|x|^3$。直接求散度为零，而$*B^\flat$在任意包围原点、外向定向的球面上的积分为$4\pi$。若有全局$\beta$使$\dd\beta=*B^\flat$，拉回球面后由无边界Stokes积分应为零，矛盾。这里不使用一个包含删点的实心球作为合法积分区域。此例说明局部势与全局势的差别，也不能用局部Hodge星替代全局上同调判断。
\end{proof}

\Needspace{12\baselineskip}
\subsection*{习题14.4：平坦环面的一次形式分解}
原创。在$\mathbb T^2$上取$g=\dd x^2+\dd y^2$，将
$\alpha=\sin x\,\dd x+\sin x\,\dd y+A\dd x+B\dd y$作Hodge分解，给出一个核外Green解。再判断共形改变度量后哪些部分仍调和。
\begin{proof}[解答]
对二形式$f\dd x\wedge\dd y$，$\delta(f\dd x\wedge\dd y)=f_y\dd x-f_x\dd y$，故
\[
 h=A\dd x+B\dd y,\qquad
 e=\dd(-\cos x)=\sin x\dd x,\qquad
 c=\delta(\cos x\dd x\wedge\dd y)=\sin x\dd y.
\]
对两个系数取全环面平均得到$h$；这里$\alpha$未必闭，不能从任意代表环路的积分直接读取其调和分量。$u=\sin x(\dd x+\dd y)$满足$\Delta_Hu=\alpha-h$，系数均零平均，所以$u\perp\mathcal H^1$，即$G\alpha=u$。其中$\delta u=-\cos x$，$\dd u=\cos x\dd x\wedge\dd y$，与分解相符。

恰当、余恰当项逐点内积为零仅是此例的特殊性；一般定理只给积分正交。若$\widetilde g=e^{2v}g$，二维一次形式上的星不变，所以一次形式闭且余闭的条件仍为$\dd\alpha=\dd*\alpha=0$，调和一次形式空间仍由$\dd x,\dd y$张成。其$L^2$内积也因逆度量与体积因子抵消而不变。对余恰当项进一步核对
\[
 \delta_{\widetilde g}(e^{2v}\cos x\,\mathrm{vol}_g)
 =-*_{g}\dd(\cos x)=\sin x\,\dd y,
\]
便知$e,c$仍为这两个正交分量；二形式的势改变了，不能把不同次数的星都当作共形不变。两者非零且分别恰当、余恰当，故均不调和。
\end{proof}

\Needspace{12\baselineskip}
\subsection*{习题14.5：球面Poisson方程的相容条件}
原创。在连通闭黎曼流形上证明$\Delta_Hu=f$有光滑函数解当且仅当$\int_Mf\mathrm{vol}=0$，解只差常数。对半径$R$球面、$f=\cos\theta$给出零平均解，并解释$f=1$为何无解。
\begin{proof}[解答]
必要性由$1$在核中及自伴性：$(f,1)=(\Delta_Hu,1)=0$。函数核由能量恒等式恰为常数，故零积分意味着$f\perp\mathcal H^0$，取$u=Gf$便得光滑解和零平均。两个解之差在函数核中，因而为常数。

球面函数只依赖$\theta$时，
\[
 \Delta_gu=\frac1{R^2\sin\theta}\partial_\theta(\sin\theta\,\partial_\theta u),
 \qquad\Delta_g\cos\theta=-\frac2{R^2}\cos\theta.
\]
所以$u=R^2\cos\theta/2$满足$\Delta_Hu=f$。它是高度函数的倍数，极点处光滑，南北反射给积分零。$f=1$有正积分，不能满足相容条件；这里若错用$\Delta_H=\Delta_g$，解的号会相反。
\end{proof}

\Needspace{12\baselineskip}
\subsection*{习题14.6：球面三角形与极帽边界}
原创。验证球面第一卦限的三条边为测地线、三角均为$\pi/2$，计算面积；再对北极帽$0\leq\theta\leq\theta_0<\pi$计算边界的带号测地曲率和Gauss--Bonnet两项。
\begin{proof}[解答]
三条边为坐标平面与球面的交，是大圆弧，所以$k_g=0$。在各坐标轴正端点，相邻两条弧的切向量沿另外两个坐标轴，故正交。角超为$3\pi/2-\pi=\pi/2$，$K=R^{-2}$，面积为$\pi R^2/2$；亦可用八个同面积卦限拼成整个球面核对。

极帽边界正向为$\phi$增加。其单位切向$T=(R\sin\theta_0)^{-1}\partial_\phi$，左法向$JT=-R^{-1}\partial_\theta$。由$\Gamma^\theta_{\phi\phi}=-\sin\theta\cos\theta$，
\[
 \nabla_TT=-\frac{\cot\theta_0}{R^2}\partial_\theta,
 \qquad k_g=\frac{\cot\theta_0}{R},\qquad
 \int_{\partial S}k_g\dd s=2\pi\cos\theta_0.
\]
面积积分给$\int_SK\mathrm{vol}=2\pi(1-\cos\theta_0)$，和为$2\pi$，对应圆盘Euler数1。超过半球时$\cot\theta_0<0$，负边界项是正确的，不应取绝对值；$\theta_0=0,\pi$不再是题设的正则边界圆。
\end{proof}

\Needspace{12\baselineskip}
\subsection*{习题14.7：共形圆环的内外边界号}
原创。在圆环$r_0\leq\sqrt{x^2+y^2}\leq R_0$上取$g=e^{2u}(\dd x^2+\dd y^2)$，$0<r_0<R_0$。推导$K$及$k_g\dd s$的共形变换，令$u=x^2+y^2$逐项验证Gauss--Bonnet。
\begin{proof}[解答]
对正向正交标架$e_1=e^{-u}\partial_x,e_2=e^{-u}\partial_y$，由Koszul或直接求Christoffel系数得联络一形式$a=-u_y\dd x+u_x\dd y$。因$\dd a=(u_{xx}+u_{yy})\dd x\wedge\dd y=-K\mathrm{vol}_g$，所以
\[
 K=-e^{-2u}(u_{xx}+u_{yy}).
\]
欧氏单位切向记$T_0$，其右侧外法向为$\nu_0$。切向角在共形正交标架下相同，故
\[
 k_g\dd s_g=k_0\dd s_0+a(T_0)\dd s_0
             =(k_0+\partial_{\nu_0}u)\dd s_0.
\]
内圈正向为顺时针，$\nu_0$朝向删去的圆盘。对$u=r^2$，曲率积分为$-4\pi(R_0^2-r_0^2)$；外圈通量及转角贡献为$2\pi+4\pi R_0^2$，内圈为$-2\pi-4\pi r_0^2$，总和为零。圆环无角点，Euler数为0。若内圈也取逆时针，就会破坏边界相消及这一等式。
\end{proof}

\Needspace{12\baselineskip}
\subsection*{习题14.8：四维Euler形式的乘积归一化}
原创。用四维Pfaffian展开计算$S^2(R_1)\times S^2(R_2)$的$e_g$及其积分。再对四维常曲率$K$的局部模型计算$e_g$，说明这些核算不能证明一般Chern--Gauss--Bonnet。
\begin{proof}[解答]
取正向乘积正交标架，前两个方向属于第一因子，后两个属于第二因子。第11章乘积联络使混合曲率全零，故
\[
 \Omega_{12}=R_1^{-2}\theta^1\wedge\theta^2,\quad
 \Omega_{34}=R_2^{-2}\theta^3\wedge\theta^4,\quad
 e_g=\frac{\mathrm{vol}_g}{4\pi^2R_1^2R_2^2}.
\]
两个球面面积分别为$4\pi R_1^2,4\pi R_2^2$，于是积分为4，半径因子全部相消。

四维常曲率模型有$\Omega_{ij}=K\theta^i\wedge\theta^j$。Pfaffian中的三项在排成$1234$次序后均为$+K^2\mathrm{vol}$，所以$e_g=3K^2\mathrm{vol}/(4\pi^2)$。特别不能把Pfaffian错当作只保留$12|34$配对的一个乘积；一般矩阵还有另两项。这些只是局部代数及特定积分。将它们推广为任意闭流形的Euler数公式，仍需本章尚未证明的整体比较定理。
\end{proof}

